\subsection{El determinante de diagonal uno: contribuciones del modo cero}\label{sec:B4T08}
Conservamos la notación de las Secciones~\ref{sec:B4T06} y~\ref{sec:B4T07}. Escribimos
$y=p^*_{\rm hijo}-w$ y
\[
 f_i=\sum_{\{a,b\}\subset a_i}\Gamma_{ab},\qquad
 f_i^I=\sum_{\{a,b\}\subset a_i\cap I}\Gamma_{ab},
 \qquad K_0=\Gamma-I.
\]
Las sumas recorren los pares no ordenados de líneas distintas.
Todas las $f_i$ son $O(k^{-1})$.
La perturbación $\Delta_0=S_{0,\rm hijo}-S_{0,II}$ tiene diagonal cero y
\[
 \operatorname{tr}(\Gamma_{II}\Delta_0)
 =-2\sum_Fp_i f_i^I+2\sum_{E'}y_i f_i.
\]
La identidad es exacta: cada par de líneas tiene a lo sumo una casilla de intersección.

\paragraph{Schur positivo y covolumen.}
Para controlar el modo cero sin una inversa de gauge mal condicionada,
ponemos $\Pi=vv^T/(2k)$ y $\widetilde S_0=S_0+\Pi$.
Su inversa es $\Gamma+\Pi$. Sean
\[
 C_J=S_{0,JJ}+\Pi_{JJ},\quad V_J=S_{0,IJ}+\Pi_{IJ},\quad
 P_*=\widetilde S_{0,II}-V_J C_J^{-1}V_J^T,
\]
\[
 H_*=P_*^{-1}=(\Gamma+\Pi)_{II},\quad
 \Pi'=v_Iv_I^T/[2(k-1)],\quad
 D_*=\Delta_0+(\Pi'-\Pi_{II})+V_JC_J^{-1}V_J^T.
\]
Las identidades exactas de Schur y cofactor son
\[
 \frac{\det S_{0,\rm hijo,gf}}{\det S_{0,\rm padre,gf}}
 =\frac{k}{k-1}\frac{\det(I+H_*D_*)}{\det C_J}.
\]
El factor $k/(k-1)$ sólo aparece al introducir estas matrices regularizadas;
es compatible con la identidad directa de cofactores de la Sección~\ref{sec:B4T01}.

\paragraph{Cuatro cálculos de primer orden.}
Como $S_{0,JJ}=I+p_c(\mathbf1\mathbf1^T-I)$,
\[
 \log[k/(k-1)]-\log\det C_J=O(k^{-2}).
\]
Las dos correcciones de proyectores satisfacen
\[
 \operatorname{tr}(\Pi_{II}\Delta_0)=-1/k+O(k^{-2}),\qquad
 \operatorname{tr}\big[H_*(\Pi'-\Pi_{II})\big]=1/k+O(k^{-2}).
\]
Se usan $\Delta_0v_I=S_{0,IJ}v_J$ y $\Gamma v=0$.
Para el término de rango cuatro, la inversa por bloques da
\[
 \operatorname{tr}(H_*V_JC_J^{-1}V_J^T)
 =\operatorname{tr}\big[C_J(\Gamma+\Pi)_{JJ}-I\big]
 =\sum_{j\in J}(K_0)_{jj}+2/k+O(k^{-2}).
\]

\paragraph{Contribución del término cuadrático.}
La parte principal de $\Delta_0$ es la matriz dispersa con entrada $-p_i$
en cada par de las tres líneas retenidas de una casilla $i\in F\setminus c$.
Su norma $\infty$ es $O(k^{-1})$. Las demás entradas de $D_*$ son
$O(k^{-2})$, con suma absoluta por fila $O(k^{-1})$.
Como $H_*=I+O_{\max}(k^{-1})$ y $\|H_*\|_\infty=O(1)$,
\[
 \operatorname{tr}(H_*D_*H_*D_*)=6\sum_Fp_i^2+O(k^{-2}).
\]
No se superponen pares de líneas de casillas distintas. La contribución
de $c$ es sólo $O(k^{-2})$. Los términos cruzados con la parte densa son
$O(k^{-2})$, y el resto cúbico del logaritmo del determinante también lo es.
Reuniendo estos cálculos,
\begin{equation}\label{eq:det-S0}
 \boxed{\log\frac{\det S_{0,\rm hijo,gf}}{\det S_{0,\rm padre,gf}}
 =\sum_{j\in J}(K_0)_{jj}+\frac2k+
 \operatorname{tr}(\Gamma_{II}\Delta_0)-3\sum_Fp_i^2+O(k^{-2}).}
\end{equation}

\paragraph{Identidad para la cancelación de primer orden.}
Sumar las cuatro entradas diagonales de $S_0\Gamma=I-\Pi$ da
\begin{equation}\label{eq:identidad-F}
 \boxed{\sum_Fp_i(f_i-f_i^I)
 =-\sum_{j\in J}(K_0)_{jj}-\frac1k-p_cf_c.}
\end{equation}
Cada casilla de $F\setminus c$ toca exactamente una línea de $J$.
La casilla $c$ cuenta cada par dos veces en la resolvente en vez de una;
la diferencia exacta es $p_cf_c=O(k^{-2})$.
Esta identidad utiliza el modo cero exacto, no una independencia supuesta.

Reunimos lo probado en esta sección.
\begin{lemasub}[Modo cero]\label{lem:modocero}
En las hipótesis de la Sección~\ref{sec:B4T01} y con la notación de esta sección, valen
\eqref{eq:det-S0} y la identidad exacta~\eqref{eq:identidad-F}.
\end{lemasub}
